\documentclass[11pt]{article}
\usepackage[T1]{fontenc}\usepackage{amsmath,amssymb,amsthm,lmodern,geometry}
\geometry{margin=.95in}\setlength{\emergencystretch}{2em}
\newtheorem{theorem}{Theorem}\newtheorem{lemma}{Lemma}
\newcommand{\dd}{\,d}\newcommand{\E}{\mathbb E}
\title{Gate threshold and a rank obstruction to sustained exponential escape}
\author{Standalone analytical increment}\date{Version 1, November 2026}
\begin{document}\maketitle
\begin{abstract}
For the unclamped critical right-half shape, an explicit monotone
integral determines a unique initial gate threshold at saturation.
The threshold increases strictly with the probe offset. We give the
subunit-mass phase diagram and the distinct sufficient escape test of RPS.
At saturation a receiver with upper-tail fraction \(f>0\) has a
strictly smaller growth ceiling than an isolated receiver. For the
continuously distributed upper tail this also obstructs a sustained
positive exponential lower rate for a fixed upper segment: two interior
ranks in such a segment would separate faster than the higher rank
can grow. The obstruction concerns exponential escape on diverging
time intervals. It allows finite-time crossings and gives no eventual
capture theorem.
\end{abstract}

\section{Scope and exact entry threshold}\label{an04goev1:scope}
Fix \(13/20\le\rho\le7/10,\ \lambda=\rho^2,\ \omega=(1-\lambda)/2\),
\(a=\rho K(\rho a),\ K(t)=\phi(t)+t\Phi(t)\). Use the CFS critical
family \(S_0=cq^{1/\omega}\), \(b=(c/A)^\omega>0\), fixed positive
\(c,A\), offsets \(\zeta>a\) on compact subsets of this domain.
The CFS/RPS entry laws concern the positive-noise target-only
comparison at \(T_A=\log(A/S_0)\). Here their limiting right-half
shape, normalized to a prescribed \(0<M_0\le1\), specifies an
initial state for the defined leading system. Its Q4 half is an atom
and its Q1 half is continuous. Q2 and Q3 do not enter this state.

WIT popsystem/P leading give the characteristic equations for positive
finite measures on chart-valid intervals with sufficient moments.
RPS moments, order and cone apply in one field with \(N=0\),
\(0<M\le1\), originally on compact support. We use the initial
integrals on the unclamped law and establish the estimates below on
any interval of its characteristic evolution with locally finite first
absolute moments. The long-time conclusion is a necessary obstruction
on such intervals; it requires no assumption that a finite-noise flow
tracks this evolution. The positive-noise probe transfer on the deep
windows, the coherent orbit, and the P resident linearization have
different hypotheses and are not used as escape estimates.

Define the CFS increasing map, on \(w>a\), by
\[
 B(w)=[w-\rho K(\rho w)]/2,\quad
 H(w)=w\exp\left[-\omega\int_w^\infty
       (B(r)^{-1}-(\omega r)^{-1})\,dr\right],\qquad L=H^{-1}.
\]
CFS gives \(H(a+)=0,\ H(w)/w\to1,\ H'/H=\omega/B\). Thus
\(H'\to1\), \(L(0+)=a\), \(L(h)\le C(1+h)\). These bounds are
uniform on the ratio box by the CFS Gaussian tail and simple-root
estimates. Write \(X_b\) for the random input under unit mass. For
any integrable test function \(j\),
\begin{equation}\label{an04goev1:integration}
 \langle j\rangle_b=\tfrac12 j(a)+\frac2\pi\int_0^\infty
                    \frac{j(L(bu))}{(1+u^2)^2}\dd u .
\end{equation}
Initially \(y=x\). The raw CFS Q1 density is asymptotic to
\(Ab^3/(2\pi x^4)\) as \(x\to\infty\); indeed it is
\(A H'(x)[1+(H(x)/b)^2]^{-2}/(2\pi b)\). Upon normalizing by
\(A/4\) its coefficient is \(2b^3/\pi\). There is no singular mass
loss at \(a\) by \eqref{an04goev1:integration}. Both first and second
absolute moments are finite (the third diverges). In particular
\(\E X_b/b\to1/\pi\) as \(b\to\infty\) by dominated
convergence in \eqref{an04goev1:integration}.

Put \(t=\rho\zeta\),
\[
 F(t,r)=\E[(Z+r)_+1_{Z>-t}]
       =\phi(\min(t,r))+r\Phi(\min(t,r)),\quad
 S_b(\zeta)=\langle F(t,\rho X_b)\rangle_b .
\]
The actual donor is \(M_0 S_b\).
\begin{theorem}\label{an04goev1:threshold}
At \(M_0=1\) there is a unique \(b_c>0\) solving
\(S_{b_c}(\zeta)=K(t)\). Initial input velocity at the aligned gate
is positive below \(b_c\), zero there, negative above. The function
\(b_c\) is continuous in \((\rho,\zeta)\) and strictly increasing in
\(\zeta\) at fixed \(\rho\).

For \(0<M_0\le1\), set
\[
 m_{\rm gate}=\frac{\zeta-\rho K(t)}{\zeta-a}.
\]
If \(M_0\le m_{\rm gate}\) every \(b>0\) gives inward velocity.
Otherwise there is a unique threshold, with the same three signs,
given by
\[
 S_b(\zeta)=\frac{\rho K(t)-(1-M_0)\zeta}{\rho M_0}.
\]
It increases with \(M_0\) on this existence range.
\end{theorem}
\begin{proof}
The velocity doubled is \(\rho K(t)-(1-M_0)\zeta-\rho M_0 S_b\).
The overlap is strictly increasing in its second argument since
\(F_r=\Phi(\min(t,r))>0\). The continuously distributed half
in \eqref{an04goev1:integration} moves strictly with \(b\).
The linear bound on \(L\) gives continuity and the limit
\(S_{0+}=K(\rho a)=a/\rho\) by domination. On \(r>t\), \(F(t,r)\)
grows linearly with slope \(\Phi(t)>0\); hence \(S_b\to\infty\).
The comparison \(\rho K(t)<\zeta\) follows from the resident root
and slope \(\lambda\Phi<1\), and \(\rho K(t)>a\).
These observations give the threshold set including its equality
case at \(M_0=m_{\rm gate}\). The right side of the threshold
equation increases strictly in \(M_0\).

For dependence on offset, write \(r=\rho X_b\), holding its law fixed.
At a zero of \(K(t)-\E F(t,r)\),
\[
 \Phi(t)\E(r-t)_+=\E[(K(t)-K(r))1_{r<t}]\le K(t).
\]
Since \(F_t=(r-t)_+\phi(t)\), differentiation under the finite first
moment yields at this zero
\[
 \partial_t(K(t)-\E F(t,r))\ge
 \Phi(t)-\phi(t)K(t)/\Phi(t)>0 .
\]
The last strict sign is \(\Phi(t)\) times the variance of a standard
normal conditioned on \(Z>-t\), from its first two truncated moments.
Thus every offset zero at fixed \(b\) is an upward crossing; two such
zeros on one interval are impossible. This proves the strict increase
of \(b_c\). Continuity follows by bracketing each root with two
strict signs and joint continuity of the integrals (including the
CFS inverse on compact parameter sets).
\end{proof}

For RPS cone the initial \(g\)-test uses \(Y=M_0\E X_b\)
and \(\rho K_s=M_0\langle\rho K(\rho X_b)\rangle_b\le Y\) since
\(X_b\ge a\). Precisely, \(g>0\) iff
\[
 M_0>m_g=\frac{\zeta+1-\rho K(t)}{\zeta+1-a},\qquad
 0<b<b_g,\qquad
 \E X_{b_g}
 =\frac{\rho K(t)-(1-M_0)(\zeta+1)}{M_0}.
\]
Here the notation on the left means the first moment from
\eqref{an04goev1:integration} at scale \(b_g\).
Existence and uniqueness follow from the same strict monotonicity.
At or below \(m_g\) the sufficient set is empty. Where \(b_g\)
exists it is strictly smaller than the exact gate threshold, since
\(\rho F(t,\rho x)\le\rho K(\rho x)<x\) for \(x>a\).
These tests concern initial gate direction and a sufficient
finite-time barrier. Neither one asserts a long-time exponential law.

\section{Donor costs at saturation}\label{an04goev1:large}
Now \(M=1,N=0\) and normalized ranks carry fixed mass \(d\sigma\).
With current donors \(S(x)=\int F(\rho x,\rho\tilde x)d\mu\),
\(K_s=\int K(\rho\tilde x)d\mu,\ Y=\int\tilde y\,d\mu\), write
\(U=x'\), \(V=y'\) for the strong receiver velocities only. Exactly,
for \(t=\rho x\ge0\),
\begin{align*}
 2U &=\lambda y-\rho K_s+
       \rho\phi(t)-\lambda y\Phi(-t)+\rho(K_s-S(x)),\\
 2V &=\lambda x-Y+\rho K(-t),\qquad \text{for }x\ge0,\\
 0&\le K_s-S(x)\le \rho\Phi(-t)\int\tilde x_+\,d\mu .
\end{align*}
The last bound integrates the discarded gate \(Z\le-t\le0\).
Thus the receiver linear part is off-diagonal \(\lambda/2\), with
the full common forces \(-\rho K_s/2,-Y/2\). Linear donor dependence
belongs in those integrals. In particular
\(0\le K_s-\rho\int\tilde x_+\,d\mu\le\phi(0)\).
One cannot replace common forces by bounded remainders during escape.
The displayed Gaussian errors are absolute bounds at each time,
including their moment and \(y\) factors.

\begin{lemma}[Per-rank growth ceiling]\label{an04goev1:ceiling}
Along a co-ordered characteristic of upper mass at least \(f>0\)
(including its own equality plateau), with the specified entry,
\[
 \max(x_+,y_+)\le C_f e^{r_+(f)s},\qquad
 r_+(f)=\frac{\lambda}{4}(\sqrt{4+f^2}-f)<\frac{\lambda}{2}
\]
for elapsed leading time \(s\ge0\) in its interval of existence.
\(C_f\) may depend on the initial coordinates. The isolated
upper comparison \(f=0\) has exponent \(\lambda/2\).
\end{lemma}
\begin{proof}
RPS's moment equation \(2Y'=-Y+\rho K_s\) and positive initial
\(Y\) imply \(Y\ge0\). For each of the upper labels
\(F(\rho x,\rho\tilde x)\ge K(\rho x)\). Hence
\(S(x)\ge fK(\rho x)\) using the actual upper mass.
At \(x\ge0\), \(K(\rho x)\ge\rho x\).
The two positive parts obey the cooperative upper comparison
\[
 2D^+\binom{x_+}{y_+}\le
 \lambda\begin{pmatrix}-f&1\\1&0\end{pmatrix}
          \binom{x_+}{y_+}+\rho\phi(0)\binom11 .
\]
At a zero face the off-diagonal and constant terms are nonnegative;
for a negative coordinate its positive part is locally constant.
For the output estimate use \(K(\rho x)\le\rho x_++\phi(0)\).
A positive left eigenvector for the largest eigenvalue reduces this
inequality to a scalar Gronwall estimate with exponent \(r_+(f)>0\).
This is an envelope rate, not an instantaneous bound on \(x'/x\)
for an arbitrary aspect ratio.
\end{proof}

For comparison, the isolated rate is attained by a passive aligned
receiver in the prescribed coherent field \(\mu=\delta_{(a,a)}\).
On \(x=y>a\), \(S(x)=K(\rho a)\), \(\rho K(\rho a)=a\), and
\[
  x'=y'=\tfrac12[\rho K(\rho x)-a]>0 .
\]
It reaches any fixed larger level in finite time. Beyond \(a/\lambda+1\),
\(x'\ge\lambda(x-a/\lambda)/2\), while the \(f=0\) upper comparison
bounds growth by \(C e^{\lambda s/2}\). Thus
\(T_R^{\rm passive}=(2/\lambda)\log R+O(1)\) for this specific
zero-mass receiver. Positive mass introduces a different donor field.

\section{What escape can be certified}\label{an04goev1:escape}
We first pay the moment costs without a cone assumption. Put
\(D(s)=1+\int(|x|+|y|)d\mu\). From the full equations,
\(S(x)\le K_s\le\phi(0)+\rho\int x_+d\mu\) and \(|Y|\le D-1\).
Integrating the absolute velocity inequalities gives \(D^+D\le7D\)
(the bound \(|U|+|V|\le3(|x|+|y|+D)\) suffices).
Consequently \(D(s)\le D(0)e^{7s}\). The weak moment identity is
justified by truncation of the test functions and finite first moments.
For two finite ranks local integrability of these moments bounds each
receiver on finite intervals by Gronwall. The receiver Jacobian on
their segment is locally bounded and its off-diagonals are positive,
so CFS's pairwise order argument holds also for the unclamped donor
law. In particular ranks in the continuous Q1 part with distinct
initial \(x\) keep a strictly positive pair gap on finite intervals.
All moment estimates include these upper labels.

\paragraph{Finite horizon lower rate.}
Fix \(h>0\), \(C_h=2D(0)e^{7h}\), and let \(\xi_f\) be the initial
\(1-f\) quantile, \(0<f<1/2\), computed from
\eqref{an04goev1:integration}. On the nonempty range
\(\xi_f>C_h/\lambda\), define
\[
 r_-(f,h)=\tfrac12(\lambda-C_h/\xi_f)>0,\qquad
 w(s)=\xi_f e^{r_-(f,h)s}.
\]
Then the entire upper mass \(f\) stays in \(x,y\ge w(s)\) on
\([0,h]\), or to the endpoint of the leading interval if earlier.
The range is nonempty because \(\xi_f\to\infty\) as \(f\downarrow0\).
At the face \(x=w,\ y\ge w\),
\(2U\ge\rho K(\rho w)-\rho K_s\). At the face \(y=w,\ x\ge w\),
\(2V\ge\rho K(\rho w)-Y\). Both are strictly greater than
\(\lambda w-C_h\ge2w'\), by the bound on \(D\) and the strict
Gaussian inequality \(K(t)>t\). The entry face and corner satisfy
the inequalities by alignment, and the terminal face ends the
estimate. Thus the certified radius at time \(h\) is at least
\(\xi_f e^{r_-(f,h)h}\). The estimate gives a positive fraction for
every fixed horizon, with the fraction chosen for that horizon.
It gives no fixed \(f,r_->0\) on diverging horizons.

\begin{theorem}[Obstruction for the unclamped continuous upper tail]
\label{an04goev1:gap}
For the specified \(q\)-independent entry law at \(M=1,\ N=0\),
no upper segment of fixed mass \(f>0\) can obey a uniform bound
\(x(s)\ge c_0 e^{rs}\), \(c_0,r>0\), from a fixed time onward on
intervals of unbounded length. This already fails at the level of
the defined equations, even if no lower bound on \(y\) is demanded.
\end{theorem}
\begin{proof}
Choose two fixed interior ranks \(i<j<1\) in the continuous portion
of the proposed segment. Put \(h_x=x_j-x_i,\ h_y=y_j-y_i\).
Both are nonnegative. Their sum initially is positive, and remains so.
Lemma~\ref{an04goev1:ceiling} gives \(x_j^+\le C_j e^{r_js}\),
\(r_j=r_+(1-j)<\lambda/2\).
For the absolute coordinates of the two labels a coarser
\(C_{ij}e^{8s}\) bound follows by the preceding \(D\)-estimate and
the pointwise linear-growth inequality (or by scalar Gronwall).
This suffices to compare polynomial factors to the Gaussian below.

Along the segment between the two states the receiver Jacobian is exactly
\[
 2J=\begin{pmatrix}
 \rho^2\phi(t)\big[\rho y-t-\int(\rho\tilde x-t)_+d\mu\big]
       &\lambda\Phi(t)\\
 \lambda\Phi(t)&0
 \end{pmatrix},\qquad t=\rho x .
\]
Indeed \(F_t(t,u)=(u-t)_+\phi(t)\). Differentiate at fixed
donors, which is legitimate with a finite first moment. Subtraction
of two velocities in this field cancels the common output mean.
Under the proposed lower bound the entire joining segment has
\(t\ge\rho c_0 e^{rs}\). Its diagonal error is bounded by
\(C_{ij}e^{8s}\phi(\rho c_0e^{rs})=o(1)\);
its off-diagonals converge to \(\lambda/2\). This step charges
the outliers in the moment integral before using the Gaussian tail.

Choose \(\epsilon>0\) so small that \(r_j<\lambda/2-2\epsilon\)
and \(l=\lambda/2-\epsilon>0\).
For all sufficiently late times on the candidate interval the
averaged Jacobian has first diagonal at least \(-\epsilon\)
and both off-diagonals at least \(l\). With \(d=h_x+h_y\) this gives
\[
 d'\ge(\lambda/2-2\epsilon)d,\qquad
 h_x'\ge l d-(l+\epsilon)h_x .
\]
The first inequality and a last-unit integration of the second give
\(h_x(s)\ge c_{ij}\exp[(\lambda/2-2\epsilon)s]\) for all later
times. The constants can be fixed before the terminal time:
on the finite initial interval the diagonal is bounded below using
\(D(0)e^{7s}\) and the two initial coordinates, hence \(d\)
has a positive lower bound at the start of this comparison.
But \(0\le h_x\le x_j^+\) once \(x_i>0\).
The necessary inequality
\begin{equation}\label{an04goev1:failure}
 c_{ij}e^{(\lambda/2-2\epsilon)s}\le C_j e^{r_js}
\end{equation}
fails on arbitrarily long intervals.
\end{proof}
The two attempts have different outcomes. The explicit orthant
works on a restricted finite horizon: its unpaid lower faces for
a proposed constant rate \(r\) are
\(\rho K(\rho w)-\rho K_s\ge2rw\) and
\(\rho K(\rho w)-Y\ge2rw\) for the sufficient moment test.
Mass fraction alone supplies no upper bound for these two donor
moments. More decisively, \eqref{an04goev1:failure} rules out the
desired sustained lower envelope for the continuous upper segment.
It makes no claim about a single passive label, the very last
label of a compact truncation, or a top atom which has no internal
pair of distinct ranks. It also does not preclude slower or
intermittent motion. A strict decrease in a per-rank \emph{ceiling}
is proved; identifying that ceiling with a realized growth rate
would give a different and unsupported statement.

\section{Radius, time comparisons and handoff}\label{an04goev1:clock}
The leading equations are mathematically defined for finite
coordinates (and sufficient donor moments). For interpreting an
expansion about the strong axis a radius \(R_q\) must at least
satisfy \(qR_q=o(1)\) for both scaled angles. There is no largest
such radius: \(\sqrt{R_q/q}\) is larger if \(qR_q\to0\) and still
has that property. Neither RPS nor the present estimates prove
uniform original-to-leading error on an arbitrary growing chart.
At \(x\asymp1/q\) the angle is order one. The unclamped law already
has mass \(O(R_q^{-3})\) outside \(R_q\) at entry.
An edge time must refer to specified finite starting ranks.

Write \(\ell=\log(1/q)\), take a fixed finite initial radius \(x_0>0\),
and set \(\log R_q=\beta\ell+o(\ell)\), \(0<\beta\le1\).
For \(\beta=1\) an inner example is \(R_q=1/[q\log(1/q)]\).
With a genuine two-sided exponential law of rate \(r>0\) up to
first hit one would have \(T_R=(\beta/r)\ell+o(\ell)\).
A lower exponential barrier only gives
\(T_R\le r_-^{-1}\log(R_q/x_0)\) if it persists that long.
An upper bound gives a lower bound on first-hit time (infinity
allowed). In particular a fixed rank with upper mass \(f\) has
\[
 T_R\ \ge\ \frac{\log R_q-O_f(1)}{r_+(f)} .
\]
For the candidate logistic weak seed \(N_0\asymp
S_0^{1-\lambda}=c^{1-\lambda}q^2\), its candidate delay to
half mass is \(\lambda^{-1}\log[(1-N_0)/N_0]
=(2/\lambda)\ell+O(1)\). Both the seed and its application across
the switch are unproved. Comparing coefficients:
\begin{itemize}
\item An isolated envelope at \(\beta=1\) has earliest coefficient
\(2/\lambda\). Equality of actual times does not follow.
For every fixed \(f>0\) and \(\beta=1\), the ceiling gives a strictly
later earliest coefficient. This also applies to a hypothetical
continuation of the zero-weak-mass system to the outer edge; the
outer expansion itself is unpaid.
\item More generally a fixed rank cannot hit before the candidate
delay to leading order if \(r_+(f)\le\beta\lambda/2\), equivalently
\(f\ge(1-\beta^2)/\beta\), when that condition is feasible.
Strict inequality in the rate comparison gives strictly later or
no hit. Otherwise the ceiling alone gives no ordering.
\item For a hypothetical two-sided rate \(r\), before, same leading
scale, and after correspond respectively to \(r>\beta\lambda/2\),
\(r=\beta\lambda/2\), and \(r<\beta\lambda/2\).
Theorem~\ref{an04goev1:gap} prevents using a fixed positive
upper-segment lower rate to predict such a clock for this law.
\end{itemize}
No weak evolution has been run in making these clock comparisons.
S7c/S8 need an outer angular field and controlled inner-to-outer
matching, with the mass and moments of the ranks approaching the
edge. They also need the actual weak entry clock, seed and coupled
strong/weak evolution, including any escaped strong-ancestry contribution
near the weak axis. Probe exposure, compensation, exact signed lag and
rate-weighted producers must be controlled in that changing field.
An initial gate sign and a radius-versus-time ceiling do not supply
those data.
\end{document}
